\documentclass[a4paper,12pt]{article}
\usepackage[utf8]{inputenc}
\usepackage[T1]{fontenc}
\usepackage[ngerman]{babel}
\usepackage{amsmath, amssymb}
\usepackage{geometry}
\usepackage{parskip}
\usepackage{titlesec}
\usepackage{enumitem}
\usepackage{array}
\usepackage{booktabs}

\geometry{left=3cm, right=3cm, top=2.5cm, bottom=2.5cm}

\titleformat{\section}{\large\bfseries}{}{0em}{}[\titlerule]
\titleformat{\subsection}{\normalsize\bfseries\scshape}{}{0em}{}

% Glossar-Eintrag: Nummer, fetter Titel, Inhalt
\newcommand{\gentry}[3]{%
  \noindent\textbf{#1\quad #2}\nopagebreak\\[0.2em]
  #3\par\smallskip
}

\begin{document}

\begin{center}
  {\LARGE\bfseries Glossar zu Lektion 5}\\[0.5em]
  {\large Mathematische Grundlagen (Modul 61111)}
\end{center}

\vspace{1em}

Dieses Glossar fasst die wesentlichen Definitionen, Merkregeln und Sätze der
fünften Lektion kompakt zusammen. Nummerierungen entsprechen dem Lehrtext.

% ======================================================================
\section*{Kapitel 14 \quad Funktionen}

\subsection*{14.1 \quad Grundlagen und erste Beispiele}

\gentry{14.1.1}{Funktion}{%
  Sei $D \subseteq \mathbb{R}$ nicht leer. Eine \textbf{Funktion} $f$ auf $D$ ist eine
  Abbildung $f: D \to \mathbb{R}$. Das \textbf{Bild} ist
  $f(D) = \{y \in \mathbb{R} \mid \exists\, x \in D: f(x) = y\}$.
}

\gentry{14.1.2}{Graph}{%
  Der \textbf{Graph} von $f: D \to \mathbb{R}$ ist die Menge
  $\{(x, f(x)) \mid x \in D\} \subseteq D \times \mathbb{R}$.
}

\gentry{14.1.3}{Dirichletfunktion}{%
  $f: \mathbb{R} \to \mathbb{R}$, definiert durch $f(x) = 1$ für $x \in \mathbb{Q}$ und
  $f(x) = 0$ für $x \in \mathbb{R} \setminus \mathbb{Q}$. Nicht grafisch darstellbar;
  in keinem Punkt stetig.
}

\gentry{14.1.4}{Komposition (Analysis)}{%
  Seien $f: D \to \mathbb{R}$ und $g: D' \to \mathbb{R}$ mit $f(D) \subseteq D'$.
  Dann ist $(g \circ f)(x) = g(f(x))$ für alle $x \in D$.
  \textit{Unterschied zur LA:} Es wird nur $f(D) \subseteq D'$ gefordert, nicht
  $f(D) = D'$.
}

\gentry{14.1.6}{Bemerkung (Umkehrfunktion)}{%
  Eine Funktion $f: D \to \mathbb{R}$ besitzt genau dann eine Umkehrfunktion
  $f^{-1}: f(D) \to \mathbb{R}$, wenn $f$ \textbf{injektiv} ist.
}

\gentry{14.1.7}{Umkehrfunktion}{%
  Sei $f: D \to \mathbb{R}$ injektiv. Die Funktion
  $f^{-1}: f(D) \to \mathbb{R}$ mit $f^{-1} \circ f = \mathrm{id}_D$ und
  $f \circ f^{-1} = \mathrm{id}_{f(D)}$ heißt \textbf{Umkehrfunktion} von $f$.
}

\gentry{14.1.9}{Monotonie}{%
  $f: D \to \mathbb{R}$ heißt \textbf{monoton wachsend} (bzw.\ \textbf{fallend}),
  falls $x < y$ $\Rightarrow$ $f(x) \le f(y)$ (bzw.\ $f(x) \ge f(y)$).
  Sie heißt \textbf{streng monoton wachsend} (bzw.\ \textbf{fallend}),
  falls $x < y$ $\Rightarrow$ $f(x) < f(y)$ (bzw.\ $f(x) > f(y)$).
}

\gentry{14.1.11}{Proposition (streng monoton $\Rightarrow$ injektiv)}{%
  Jede streng monotone Funktion $f: D \to \mathbb{R}$ ist injektiv.
  Ist $f$ streng monoton wachsend (bzw.\ fallend), so ist auch $f^{-1}$ streng
  monoton wachsend (bzw.\ fallend).
}

\gentry{14.1.14}{Merkregel (Graph der Umkehrfunktion)}{%
  Den Graph von $f^{-1}$ erhält man durch Spiegelung des Graphen von $f$
  an der Diagonale $d = \{(z, z) \mid z \in \mathbb{R}\}$.
}

\gentry{14.1.15}{Algebraische Verknüpfungen von Funktionen}{%
  Seien $f, g: D \to \mathbb{R}$ und $a \in \mathbb{R}$. Definiert für alle $x \in D$:
  $(f \pm g)(x) = f(x) \pm g(x)$;\; $(af)(x) = af(x)$;\; $(fg)(x) = f(x)g(x)$.
  Falls $g(x) \neq 0$: $\left(\frac{f}{g}\right)(x) = \frac{f(x)}{g(x)}$.
}

\gentry{14.1.19}{Restriktion}{%
  Sei $D' \subseteq D$, $D' \neq \emptyset$. Die Funktion
  $f|_{D'}: D' \to \mathbb{R}$, $x \mapsto f(x)$, heißt \textbf{Restriktion}
  (Einschränkung) von $f$ auf $D'$.
}

\gentry{14.1.20}{Beschränktheit}{%
  $f: D \to \mathbb{R}$ heißt \textbf{nach oben beschränkt}, wenn $\exists\, b \in \mathbb{R}$:
  $f(x) \le b$ für alle $x \in D$; \textbf{nach unten beschränkt}, wenn
  $\exists\, a \in \mathbb{R}$: $a \le f(x)$ für alle $x$. \textbf{Beschränkt}
  bedeutet beides.
}

% ======================================================================
\subsection*{14.2 \quad Beispiele}

\gentry{14.2.2}{Produkt von Polynomen}{%
  Seien $p = \sum_{i=0}^n a_i T^i$ und $q = \sum_{j=0}^m b_j T^j \in K[T]$.
  Das Produkt ist $pq = \sum_k c_k T^k$ mit
  $c_k = \sum_{i+j=k} a_i b_j$.
}

\gentry{14.2.5}{Gradformel}{%
  Seien $p, q \in K[T]$ mit $\mathrm{Grad}(p) = n \ge 0$ und $\mathrm{Grad}(q) = m \ge 0$.
  Dann gilt $\mathrm{Grad}(pq) = n + m$.
}

\gentry{14.2.6}{Proposition (Division mit Rest)}{%
  Zu $p, q \in K[T]$ mit $q \neq 0$ gibt es eindeutig bestimmte $f, r \in K[T]$
  mit $p = fq + r$ und $\mathrm{Grad}(r) < \mathrm{Grad}(q)$.
}

\gentry{14.2.7}{Nullstelle}{%
  Sei $p \in K[T]$. Ein $\lambda \in K$ mit $\tilde{p}(\lambda) = \sum_i a_i \lambda^i = 0$
  heißt \textbf{Nullstelle} von $p$.
}

\gentry{14.2.8}{Proposition (Nullstellen und Faktoren)}{%
  $\lambda \in K$ ist genau dann Nullstelle von $p \in K[T]$, wenn
  $p = (T - \lambda) \cdot q$ für ein Polynom $q \in K[T]$.
}

\gentry{14.2.9}{Korollar (Anzahl der Nullstellen)}{%
  Ein Polynom $p \in K[T]$, $p \neq 0$, vom Grad $n \ge 0$ hat höchstens $n$
  verschiedene Nullstellen.
}

\gentry{14.2.10}{Polynomfunktion}{%
  Zu $p = \sum_{i=0}^n a_i T^i \in \mathbb{R}[T]$ ist die \textbf{Polynomfunktion}
  $\tilde{p}: \mathbb{R} \to \mathbb{R}$ definiert durch
  $\tilde{p}(x) = \sum_{i=0}^n a_i x^i$.
}

\gentry{14.2.13}{Satz (Identitätssatz für Polynomfunktionen)}{%
  Seien $p, q \in \mathbb{R}[T]$ mit $\mathrm{Grad}(p) = n$, $\mathrm{Grad}(q) = m$.
  Gilt $\tilde{p}(x) = \tilde{q}(x)$ für mehr als $\max(n, m)$ verschiedene $x \in \mathbb{R}$,
  so ist $m = n$ und $a_i = b_i$ für alle $i$.
}

\gentry{14.2.16}{Rationale Funktion}{%
  Seien $\tilde{p}, \tilde{q}: \mathbb{R} \to \mathbb{R}$ Polynomfunktionen. Dann heißt
  $\dfrac{\tilde{p}}{\tilde{q}}: \mathbb{R} \setminus \{\text{Nullstellen von } q\} \to \mathbb{R}$
  eine \textbf{rationale Funktion}.
}

\gentry{14.2.21}{Allgemeine Potenz $a^\rho$}{%
  Sei $a > 0$, $\rho \in \mathbb{R}$. Wähle eine Folge $(r_n)$ rationaler Zahlen mit
  $r_n \to \rho$. Dann ist $a^\rho := \lim_{n \to \infty} a^{r_n}$.
  (Der Grenzwert ist unabhängig von der Wahl von $(r_n)$.)
  Für $a = 0$, $\rho > 0$: $0^\rho = 0$.
}

\gentry{14.2.32}{Satz (Potenzregeln für reelle Exponenten)}{%
  Seien $a, b > 0$ und $\rho, \sigma \in \mathbb{R}$:
  \begin{enumerate}[label=\arabic*.,topsep=2pt,itemsep=1pt]
    \item $a^\rho a^\sigma = a^{\rho+\sigma}$
    \item $\frac{a^\rho}{a^\sigma} = a^{\rho-\sigma}$
    \item $a^\rho b^\rho = (ab)^\rho$
    \item $\frac{a^\rho}{b^\rho} = \left(\frac{a}{b}\right)^\rho$
    \item $(a^\rho)^\sigma = a^{\rho\sigma}$
  \end{enumerate}
}

\gentry{14.2.33}{Satz (Potenzfunktion)}{%
  Sei $\rho \in \mathbb{R}$, $\rho \neq 0$, $f: (0,\infty) \to \mathbb{R}$, $x \mapsto x^\rho$.
  Dann: $\mathrm{Bild}(f) = (0,\infty)$; $f$ ist streng monoton wachsend für $\rho > 0$,
  streng monoton fallend für $\rho < 0$; $f^{-1}(x) = x^{1/\rho}$.
}

\gentry{14.2.34}{Exponentialfunktion}{%
  Sei $a > 0$. Die Funktion $\mathrm{exp}_a: \mathbb{R} \to \mathbb{R}$,
  $\mathrm{exp}_a(x) = a^x$, heißt \textbf{Exponentialfunktion} zur Basis $a$.
}

\gentry{14.2.35}{Satz (Eigenschaften der Exponentialfunktion)}{%
  Sei $a > 0$:
  \begin{itemize}[topsep=2pt,itemsep=1pt]
    \item Für $a \neq 1$: $\mathrm{Bild}(\mathrm{exp}_a) = (0,\infty)$; für $a = 1$: $\mathrm{exp}_1 = \hat{1}$.
    \item $a > 1$: streng monoton wachsend; $0 < a < 1$: streng monoton fallend.
    \item $\mathrm{exp}_a(0) = 1$, $\mathrm{exp}_a(1) = a$.
    \item $\mathrm{exp}_a(x+y) = \mathrm{exp}_a(x) \cdot \mathrm{exp}_a(y)$ für alle $x, y \in \mathbb{R}$.
  \end{itemize}
}

\gentry{14.2.37}{Notation ($e$-Funktion)}{%
  Die Funktion $\mathrm{exp}_e$ wird mit $\mathrm{exp}$ (ohne Index) bezeichnet.
  Man schreibt $e^x$ für $\mathrm{exp}(x)$.
}

\gentry{14.2.39}{Satz (Grenzwertdarstellung von $e$)}{%
  Sei $(x_n)$ eine Nullfolge mit $x_n > -1$ und $x_n \neq 0$. Dann gilt
  $\displaystyle\lim_{n \to \infty} (1 + x_n)^{1/x_n} = e$.
}

\gentry{14.2.40}{Logarithmus}{%
  Sei $a > 0$, $a \neq 1$. Die Umkehrfunktion von $\mathrm{exp}_a$ heißt
  \textbf{Logarithmus zur Basis} $a$, bezeichnet $\log_a: (0,\infty) \to \mathbb{R}$.
}

\gentry{14.2.42}{Satz (Eigenschaften des Logarithmus)}{%
  Sei $a > 0$, $a \neq 1$:
  \begin{itemize}[topsep=2pt,itemsep=1pt]
    \item $\mathrm{Bild}(\log_a) = \mathbb{R}$.
    \item $a > 1$: streng monoton wachsend; $0 < a < 1$: streng monoton fallend.
    \item $\log_a(1) = 0$, $\log_a(a) = 1$.
    \item $\log_a(xy) = \log_a(x) + \log_a(y)$ und $\log_a\!\left(\frac{x}{y}\right) = \log_a(x) - \log_a(y)$.
    \item $\log_a(x^\rho) = \rho \log_a(x)$ für alle $\rho \in \mathbb{R}$, $x > 0$.
  \end{itemize}
}

\gentry{14.2.43}{Natürlicher Logarithmus}{%
  Die Funktion $\log_e: (0,\infty) \to \mathbb{R}$ wird mit $\ln$ bezeichnet und
  \textbf{natürlicher Logarithmus} genannt.
}

\gentry{14.2.44}{Proposition (Umrechnung zwischen $\exp$, $\exp_a$, $\log_a$)}{%
  Sei $a > 0$:
  \begin{itemize}[topsep=2pt,itemsep=1pt]
    \item Für alle $x > 0$: $x^a = \exp(a \ln x)$.
    \item Für alle $x \in \mathbb{R}$: $\mathrm{exp}_a(x) = \exp(x \ln a)$.
    \item Für $a \neq 1$ und alle $x > 0$: $\log_a(x) = \dfrac{\ln(x)}{\ln(a)}$.
  \end{itemize}
}

% ======================================================================
\section*{Kapitel 15 \quad Stetigkeit}

\subsection*{15.1 \quad Grundlagen, Beispiele, erste Eigenschaften}

\gentry{15.1.1}{Stetigkeit}{%
  $f: D \to \mathbb{R}$ heißt \textbf{stetig in} $a \in D$, wenn für jede Folge
  $(a_n)$ in $D$ mit $a_n \to a$ gilt: $f(a_n) \to f(a)$.
  $f$ heißt \textbf{stetig auf} $D$, wenn $f$ in jedem $a \in D$ stetig ist.
}

\gentry{15.1.4}{Beispiele stetiger Funktionen}{%
  Die allgemeine Potenzfunktion $x \mapsto x^\rho$, die Exponentialfunktion
  $\mathrm{exp}_a$ und der Logarithmus $\log_a$ sind auf ihren jeweiligen
  Definitionsbereichen stetig.
}

\gentry{15.1.9}{Satz ($\varepsilon$-$\delta$-Kriterium für Stetigkeit)}{%
  $f: D \to \mathbb{R}$ ist genau dann in $a \in D$ stetig, wenn gilt:
  \[
    \forall\, \varepsilon > 0 \;\exists\, \delta > 0 \;\forall\, x \in D:
    \quad |x - a| < \delta \;\Rightarrow\; |f(x) - f(a)| < \varepsilon.
  \]
}

\gentry{15.1.11}{Proposition (Rechenregeln für Stetigkeit)}{%
  Seien $f, g: D \to \mathbb{R}$ stetig in $a$, $\alpha \in \mathbb{R}$.
  Dann sind $f \pm g$, $\alpha f$, $fg$ und (falls $g \neq 0$) $f/g$ stetig in $a$.
  Außerdem: Ist $h: E \to \mathbb{R}$ mit $f(D) \subseteq E$ stetig in $f(a)$,
  so ist $h \circ f$ stetig in $a$.
}

\gentry{15.1.13}{Proposition (Stetigkeit von Polynomen und rationalen Funktionen)}{%
  Polynomfunktionen sind auf ganz $\mathbb{R}$ stetig. Rationale Funktionen sind
  auf ihrem Definitionsbereich stetig.
}

% ======================================================================
\subsection*{15.2 \quad Stetige Funktionen auf Intervallen}

\gentry{15.2.2}{Satz (Nullstellensatz von Bolzano)}{%
  Sei $f: [a,b] \to \mathbb{R}$ stetig mit $f(a) < 0 < f(b)$ (oder umgekehrt).
  Dann gibt es mindestens ein $x \in (a,b)$ mit $f(x) = 0$.
}

\gentry{15.2.3}{Nullstelle}{%
  Ein $x \in D$ mit $f(x) = 0$ heißt \textbf{Nullstelle} von $f: D \to \mathbb{R}$.
}

\gentry{15.2.4}{Korollar (Zwischenwertsatz von Bolzano)}{%
  Sei $f: [a,b] \to \mathbb{R}$ stetig, und sei $d$ eine Zahl zwischen $f(a)$ und $f(b)$.
  Dann gibt es ein $x_d \in [a,b]$ mit $f(x_d) = d$.
  Insbesondere ist das Bild einer stetigen Funktion auf einem Intervall wieder
  ein Intervall.
}

\gentry{15.2.8}{Satz (Stetige Bilder kompakter Intervalle)}{%
  Sei $f: [a,b] \to \mathbb{R}$ stetig. Dann ist $f([a,b])$ ein abgeschlossenes
  und beschränktes Intervall.
}

\gentry{15.2.9}{Korollar (Satz vom Minimum und Maximum)}{%
  Sei $f: [a,b] \to \mathbb{R}$ stetig. Dann gibt es Stellen $x_1, x_2 \in [a,b]$
  mit $f(x_1) \le f(x) \le f(x_2)$ für alle $x \in [a,b]$.
}

\gentry{15.2.10}{Lemma (Stetig + injektiv $\Rightarrow$ streng monoton)}{%
  Sei $f: [a,b] \to \mathbb{R}$ injektiv und stetig. Dann ist $f$ streng monoton.
}

\gentry{15.2.12}{Korollar (Injektiv $\Leftrightarrow$ streng monoton auf Intervall)}{%
  Sei $I$ ein Intervall und $f: I \to \mathbb{R}$ stetig. Genau dann ist $f$
  injektiv, wenn $f$ streng monoton ist.
}

\gentry{15.2.14}{Proposition (Stetigkeit der Umkehrfunktion)}{%
  Sei $f: I \to \mathbb{R}$ injektiv und stetig, $I$ ein Intervall. Dann ist
  die Umkehrfunktion $f^{-1}: f(I) \to \mathbb{R}$ ebenfalls stetig.
}

% ======================================================================
\subsection*{15.3 \quad Grenzwerte von Funktionen}

\gentry{15.3.1}{Häufungspunkt}{%
  $a \in \mathbb{R}$ heißt \textbf{Häufungspunkt} von $M \subseteq \mathbb{R}$,
  wenn es eine Folge $(a_n)$ in $M \setminus \{a\}$ gibt mit $a_n \to a$.
}

\gentry{15.3.4}{Konvergenz einer Funktion}{%
  $f: D \to \mathbb{R}$ heißt \textbf{konvergent in} $a$ (einem Häufungspunkt von $D$),
  wenn $(f(a_n))$ für jede Folge $(a_n)$ in $D \setminus \{a\}$ mit $a_n \to a$
  konvergiert.
}

\gentry{15.3.8}{Definition (Grenzwert)}{%
  Ist $f$ konvergent in $a$, so heißt der Grenzwert einer solchen Folge
  $(f(a_n))$ der \textbf{Grenzwert} von $f$ in $a$:
  $\displaystyle\lim_{x \to a} f(x) = b$.
  (Dieser Grenzwert ist unabhängig von der Wahl der Folge.)
}

\gentry{15.3.9}{Proposition (Stetigkeit über Grenzwert charakterisiert)}{%
  Sei $a \in D$ ein Häufungspunkt von $D$. Dann gilt:
  $f$ ist genau dann stetig in $a$, wenn $\displaystyle\lim_{x \to a} f(x) = f(a)$.
}

\gentry{15.3.11}{Stetige Fortsetzung}{%
  Sei $a \notin D$ ein Häufungspunkt und $\lim_{x \to a} f(x) = c$.
  Die Funktion $g: D \cup \{a\} \to \mathbb{R}$ mit $g|_D = f$ und $g(a) = c$
  heißt \textbf{stetige Fortsetzung} von $f$ auf $D \cup \{a\}$.
}

\gentry{15.3.14}{Hebbare Unstetigkeit}{%
  $f: D \to \mathbb{R}$ besitzt in $a \in D$ eine \textbf{hebbare Unstetigkeit},
  wenn $f|_{D\setminus\{a\}}$ auf $D$ stetig fortgesetzt werden kann, also
  $\lim_{x \to a} f(x)$ existiert (aber $\neq f(a)$ ist).
}

\gentry{15.3.15}{Satz ($\varepsilon$-$\delta$-Kriterium für Grenzwert)}{%
  $\displaystyle\lim_{x \to a} f(x) = b$ genau dann, wenn gilt:
  \[
    \forall\, \varepsilon > 0 \;\exists\, \delta > 0 \;\forall\, x \in D \setminus \{a\}:
    \quad 0 < |x - a| < \delta \;\Rightarrow\; |f(x) - b| < \varepsilon.
  \]
}

\gentry{15.3.17}{Proposition (Rechenregeln für Funktionenkonvergenz)}{%
  Seien $\lim_{x\to a} f(x) = u$ und $\lim_{x \to a} g(x) = v$, $\alpha \in \mathbb{R}$.
  Dann gilt:
  $\lim_{x\to a}(f \pm g)(x) = u \pm v$;\;
  $\lim_{x\to a}(\alpha f)(x) = \alpha u$;\;
  $\lim_{x\to a}(fg)(x) = uv$;\;
  $\lim_{x\to a}(f/g)(x) = u/v$ (falls $v \neq 0$).
}

% ======================================================================
\section*{Kapitel 16 \quad Differenzierbarkeit}

\subsection*{16.1 \quad Die Ableitung einer Funktion}

\gentry{16.1.1}{Differenzierbarkeit und Ableitung}{%
  $f: I \to \mathbb{R}$ heißt \textbf{differenzierbar in} $a \in I$, wenn
  \[
    f'(a) := \lim_{x \to a} \frac{f(x) - f(a)}{x - a}
    = \lim_{h \to 0} \frac{f(a+h) - f(a)}{h}
  \]
  existiert. $f'(a)$ heißt \textbf{Ableitung} von $f$ in $a$.
}

\gentry{16.1.3}{Proposition (differenzierbar $\Rightarrow$ stetig)}{%
  Ist $f$ in $a$ differenzierbar, so ist $f$ in $a$ stetig.
  (Die Umkehrung gilt im Allgemeinen nicht.)
}

\gentry{16.1.5}{Differenzenquotient}{%
  Der Bruch $\dfrac{f(x) - f(a)}{x - a}$ ($x \neq a$) heißt \textbf{Differenzenquotient}.
  Er ist die Steigung der \textbf{Sekante} durch $(a, f(a))$ und $(x, f(x))$.
}

\gentry{16.1.6}{Merkregel (Tangente)}{%
  Die Ableitung $f'(a)$ ist die \textbf{Steigung der Tangente} an den Graphen von
  $f$ im Punkt $(a, f(a))$. Die Tangente hat die Gleichung $y = f(a) + f'(a)(x-a)$.
}

\gentry{16.1.7}{Lokale Extrema}{%
  $f: D \to \mathbb{R}$ hat in $a \in D$ ein \textbf{lokales Maximum} (bzw.\
  \textbf{Minimum}), wenn es eine $\delta$-Umgebung $U_\delta(a)$ gibt, sodass
  $f(x) \le f(a)$ (bzw.\ $f(x) \ge f(a)$) für alle $x \in U_\delta(a) \cap D$.
}

\gentry{16.1.8}{Globale Extrema}{%
  $f$ hat in $a$ ein \textbf{globales Maximum} (bzw.\ \textbf{Minimum}),
  wenn $f(a) \ge f(x)$ (bzw.\ $f(a) \le f(x)$) für alle $x \in D$.
  Globale Extrema sind stets auch lokale Extrema.
}

\gentry{16.1.10}{Innerer Punkt}{%
  $x \in X \subseteq \mathbb{R}$ heißt \textbf{innerer Punkt} von $X$, wenn
  $\exists\, \delta > 0: U_\delta(x) \subseteq X$.
}

\gentry{16.1.12}{Proposition (notwendige Bedingung für Extrema)}{%
  Sei $f: I \to \mathbb{R}$ in einem inneren Punkt $a \in I$ differenzierbar.
  Hat $f$ in $a$ ein lokales Extremum, so gilt $f'(a) = 0$.
  \textit{Vorsicht:} Die Umkehrung gilt nicht.
}

% ======================================================================
\subsection*{16.2 \quad Differentiationsregeln}

\gentry{16.2.1}{Proposition (Rechenregeln der Differentialrechnung)}{%
  Seien $f, g: I \to \mathbb{R}$ in $a$ differenzierbar, $\alpha \in \mathbb{R}$:
  \begin{itemize}[topsep=2pt,itemsep=1pt]
    \item \textbf{Summenregel:} $(f \pm g)'(a) = f'(a) \pm g'(a)$
    \item \textbf{Faktorregel:} $(\alpha f)'(a) = \alpha f'(a)$
    \item \textbf{Produktregel:} $(fg)'(a) = f(a)g'(a) + f'(a)g(a)$
    \item \textbf{Quotientenregel:} $\left(\frac{f}{g}\right)'(a) = \dfrac{f'(a)g(a) - f(a)g'(a)}{(g(a))^2}$
      \quad (falls $g(a) \neq 0$)
  \end{itemize}
}

\gentry{16.2.2}{Korollar (Reziprokregel)}{%
  Ist $g$ in $a$ differenzierbar mit $g(a) \neq 0$, so gilt
  $\left(\dfrac{1}{g}\right)'(a) = -\dfrac{g'(a)}{(g(a))^2}$.
}

\gentry{16.2.3}{Satz (Kettenregel)}{%
  Sei $g$ in $a$ und $f$ in $g(a)$ differenzierbar, $\mathrm{Bild}(g) \subseteq I_f$.
  Dann ist $f \circ g$ in $a$ differenzierbar, und es gilt:
  \[
    (f \circ g)'(a) = f'(g(a)) \cdot g'(a).
  \]
}

\gentry{16.2.4}{Proposition (Ableitung der Umkehrfunktion)}{%
  Sei $f: I \to \mathbb{R}$ stetig, streng monoton und in $a$ differenzierbar mit
  $f'(a) \neq 0$. Dann ist $f^{-1}$ in $b = f(a)$ differenzierbar, und es gilt:
  \[
    (f^{-1})'(b) = \frac{1}{f'(a)} = \frac{1}{f'(f^{-1}(b))}.
  \]
}

% ======================================================================
\subsection*{16.3 \quad Ableitungen wichtiger Funktionen}

\gentry{16.3.1}{Proposition (Polynomfunktionen)}{%
  Polynomfunktionen $f(x) = \sum_{i=0}^n a_i x^i$ sind auf ganz $\mathbb{R}$
  differenzierbar, und es gilt $f'(x) = \sum_{i=1}^n i a_i x^{i-1}$.
}

\gentry{16.3.2}{Proposition (ganze Potenzen)}{%
  Für alle $k \in \mathbb{Z}$ und $x \neq 0$ (falls $k < 0$) gilt:
  $(x^k)' = k x^{k-1}$.
}

\gentry{16.3.4}{Satz (Exponentialfunktion)}{%
  Für alle $a > 0$ und alle $x \in \mathbb{R}$ gilt:
  $\mathrm{exp}_a'(x) = \mathrm{exp}_a(x) \cdot \ln(a)$, insbesondere $\exp'(x) = \exp(x)$.
}

\gentry{16.3.5}{Korollar (natürlicher Logarithmus)}{%
  $\ln: (0,\infty) \to \mathbb{R}$ ist differenzierbar, und es gilt $\ln'(x) = \dfrac{1}{x}$.
}

\gentry{16.3.6}{Korollar (allgemeiner Logarithmus)}{%
  Für $a > 0$, $a \neq 1$, ist $\log_a$ differenzierbar, und es gilt
  $\log_a'(x) = \dfrac{1}{x \ln(a)}$.
}

\gentry{16.3.7}{Korollar (allgemeine Potenzfunktion)}{%
  Für $a \in \mathbb{R}$, $f(x) = x^a$ auf $(0,\infty)$: $f$ ist differenzierbar,
  und es gilt $f'(x) = a x^{a-1}$.
}

\gentry{16.3.11}{Satz (Winkelfunktionen)}{%
  Sinus und Cosinus sind auf $\mathbb{R}$ differenzierbar. Es gilt:
  \[
    \sin'(x) = \cos(x), \qquad \cos'(x) = -\sin(x).
  \]
  Verwendet werden dabei die Eigenschaften SinCos-1 bis SinCos-5,
  insbesondere $\lim_{x\to 0}\frac{\sin x}{x} = 1$.
}

\end{document}
